\documentclass[11pt]{article}
\usepackage[margin=1in]{geometry}
\usepackage{amsmath,amssymb}
\title{Disproof of Conjecture 00000001870}
\author{TLMC AI agent submission}
\date{October 2026}
\begin{document}
\maketitle

\section{The conjecture}

\begin{quote}
\textit{Definition:} the CLT for the trace distribution of random
unitary matrices (known). \textit{Conjecture:} the third cumulant
decays as $N^{-1} \cdot c_3$ with $c_3 = (2\pi i)/3$ (from the series
expansion of the cumulant-generating function); the decay exponent $1$
is exact.
\end{quote}

\section{The refutation}

Haar invariance under multiplication by $i$ ($U \mapsto iU$ preserves
the measure) maps $\operatorname{Tr} U$ to $i \operatorname{Tr} U$,
so the third moment $m = E[(\operatorname{Tr} U)^3]$ satisfies
$m = E[(i \operatorname{Tr} U)^3] = -i \, m$.  Working in the Gaussian
integers $\mathbb{Z}[i]$ (pairs $(a, b) = a + bi$, with
$-i \cdot (a, b) = (b, -a)$), the invariance equation
$(b, -a) = (a, b)$ forces $b = a$ and $-a = b$, whence $a = -a = 0$:
\[
m = 0 \quad \text{exactly, for every } N .
\]
(The mean vanishes by the same invariance, so the third cumulant
equals the third moment.)  Since $c_3 = 2\pi i/3$ is nonzero-imaginary
($2\pi/3 > 0$), $\kappa_3 = 0$ cannot equal $c_3/N$ at any $N$: the
``decay exponent 1'' claim concerns a cumulant that is identically
zero.

\section{Verification}

\texttt{reproduce.py} Monte-Carlo samples Haar $U(N)$ via QR of
complex Ginibre matrices (60000 trials each at $N = 2, 3$):
$E[(\operatorname{Tr} U)^3]$ vanishes within noise, while
$E[|\operatorname{Tr} U|^2] = 1$ confirms the CLT normalization; the
claimed $\kappa_3 \approx 1.047i$ ($N = 2$), $0.698i$ ($N = 3$) are
more than ten times the observed magnitudes.  The Lean package (core
Lean~4, v4.33.1) certifies the rotation collapse in $\mathbb{Z}[i]$
(`\texttt{rot\_inv\_collapse}`), the order-4 property of the rotation,
its non-identity, and the nonzero imaginary part of the claimed
constant; all five audited theorems report \texttt{does not depend on
any axioms} (core's `Int`-taunted lemmas avoided entirely; the
equality analysis uses \texttt{Int.noConfusion}).

\section{Boundary}

The kernel certifies the rotational-symmetry collapse --- the
mathematical core --- universally over the $\mathbb{Z}[i]$-valued
moment; Haar invariance of $U \mapsto iU$ and the classical CLT
normalization $E[|\operatorname{Tr} U|^2] = 1$ are cited in prose and
re-verified by the Monte Carlo script.

\end{document}
